% Abstract only, laid out like a NeurIPS-style paper page (single column, 5.5 in text width, Times).
\documentclass[10pt]{article}
\usepackage[letterpaper,textwidth=5.5in,textheight=9in,top=1in]{geometry}
\usepackage[T1]{fontenc}
\usepackage{mathptmx}
\usepackage{microtype}

\pagestyle{empty}
\setlength{\parindent}{0pt}

\begin{document}

\vbox{%
  \hrule height 4pt
  \vskip 0.25in
  \centering
  {\LARGE\bfseries TDP-Reasoner: Timing-Dependent Plasticity\\[2pt] Along the Recursion\par}
  \vskip 0.29in
  \hrule height 1pt
}
\vskip 0.3in
\begin{center}
  {\bfseries Anonymous Author(s)}\\
  Affiliation\\
  \texttt{email}
\end{center}
\vskip 0.3in

\begin{center}
  {\large\bfseries Abstract}
\end{center}
\vskip 0.05in
\begin{list}{}{\leftmargin=0.5in \rightmargin=0.5in}\item[]
Small recursive models trained from scratch are the strongest non-LLM approach to puzzle benchmarks such as
Sudoku-Extreme and ARC-AGI. These models, HRM and TRM among them, re-apply a Transformer block many times. Seen as a
modern Hopfield network, each step's attention is a Hebbian memory of the current state, so recursion reinforces the
model's current answer. We propose TDP-Reasoner
(timing-dependent plasticity reasoner), which keeps a fast-weight memory along the recursion itself: keys and values from
different steps are paired under a learned two-sided timing window, borrowed from spike-timing-dependent plasticity, so
the memory tracks how the state moves rather than where it is. Demonstrations of a task reach the query only through a
per-task linear rule fitted in context, so unseen tasks need no retraining. Run 39 times deeper than it was trained, it
removes 55\% of its remaining errors on held-out Sudoku-Extreme puzzles and never loses a solved puzzle. On
unseen
ARC-AGI-1 tasks with small grids it matches the public TRM checkpoint (43.2\% vs.\ 45.5\%) after about 3\% of its
training; at full scale it reaches 18.8\% pass@2 on ARC-AGI-1 (TRM: 41.8\% under the same evaluation) and 0\% on
ARC-AGI-2.
\end{list}

\end{document}
